\documentclass[12pt,a4paper]{article}

% --- Encodage et langue ---
\usepackage[utf8]{inputenc}
\usepackage[T1]{fontenc}
\usepackage[french]{babel}

% --- Police times (guide wikistat) ---
\usepackage{mathptmx}

% --- Mise en page ---
\usepackage[top=2.5cm, bottom=2.5cm, left=2.5cm, right=2.5cm]{geometry}
\usepackage{setspace}
\onehalfspacing

% --- Graphiques et tableaux ---
\usepackage{graphicx}
\usepackage{float}
\usepackage{booktabs}
\usepackage{array}
\usepackage{multirow}
\usepackage{tabularx}
\usepackage{pgfplots}
\pgfplotsset{compat=1.18}
\usepackage{tikz}
\usepackage{xurl}

% --- Maths et symboles ---
\usepackage{amsmath, amssymb}
\usepackage{eurosym}

% --- Liens ---
\usepackage[colorlinks=true, linkcolor=blue!60!black, citecolor=blue!60!black, urlcolor=blue!60!black]{hyperref}

% --- En-tetes ---
\usepackage{fancyhdr}
\pagestyle{fancy}
\fancyhf{}
\rhead{\textit{Projet RCP208}}
\lhead{\textit{March\'e immobilier d'Amiens}}
\cfoot{\thepage}

% --- Legendes figures ---
\usepackage[font=small,labelfont=bf,textfont=it]{caption}

% --- Couleurs + boites description ---
\usepackage{xcolor}
\usepackage{tcolorbox}
\tcbuselibrary{skins}

\definecolor{steelblue}{RGB}{70,130,180}

\pgfplotsset{
    rapportaxis/.style={
        tick label style={font=\small},
        label style={font=\small},
        title style={font=\small},
        grid=major,
        grid style={dashed, gray!30},
        scaled ticks=false,
    }
}

\newtcolorbox{descrbox}[1][]{%
  enhanced,
  colback=blue!3,
  colframe=blue!40!black,
  boxrule=0.5pt,
  arc=3pt,
  left=8pt, right=8pt, top=5pt, bottom=5pt,
  fontupper=\small,
  fonttitle=\small\bfseries\sffamily,
  title={\raisebox{-0.2ex}{\textsf{$\triangleright$}}~Interpr\'etation},
  attach boxed title to top left={yshift=-2mm, xshift=4mm},
  boxed title style={colback=blue!40!black, colframe=blue!40!black, arc=2pt, boxrule=0pt},
  #1
}

% --- Enumeration compacte ---
\usepackage{enumitem}
\setlist[enumerate]{topsep=2pt, itemsep=1pt, parsep=0pt}
\setlist[itemize]{topsep=2pt, itemsep=1pt, parsep=0pt}

% ============================================================
\begin{document}

% ============================================================
% COUVERTURE
% ============================================================
\begin{titlepage}
\noindent
\begin{minipage}[t]{0.3\textwidth}
\raggedright
\includegraphics[height=2cm]{cnam_logo.jpg}
\end{minipage}%
\hfill
\begin{minipage}[t]{0.3\textwidth}
\raggedleft
\includegraphics[height=3cm]{ipst_cnam_logo-removebg-preview.png}
\end{minipage}

\centering
\vspace*{1cm}
{\Large Cnam Occitanie}\\[0.3cm]
{\large IPST Cnam d'Albi}\\[0.3cm]
\textbf{{Apprentissage statistique :}} modélisation descriptive et introduction aux réseaux de neurones
\textbf{{UE : RCP208 }}\\[2.5cm]
\rule{0.8\textwidth}{1.2pt}\\[0.6cm]
{\LARGE \textbf{Projet : Estimation des prix de ventes immobilières
}}\\[0.2cm]
\rule{0.8\textwidth}{1.2pt}\\[2.5cm]
{\large Auteur : \textbf{Mohamed Amine Sobhi}}\\[0.5cm]
{\large Professeur: \textbf{ M. David Jaidan }}\\[2.5cm]
{\large Master Informatique : Parcours TRIED\par}
{\large Ann\'ee universitaire 2025-2026}\\[1.5cm]
\vfill
\begin{center}
        {\large Mai 2026\par}
    \end{center}
\end{titlepage}

% ============================================================
\newpage
\tableofcontents
\newpage

% ============================================================
% RESUME D'AUTEUR
% ============================================================
\section*{R\'esum\'e d'auteur}
\addcontentsline{toc}{section}{R\'esum\'e d'auteur}

\noindent\textit{\small Cette section est incluse conform\'ement aux recommandations du guide de r\'edaction de rapports  (Wikistat), qui la qualifie d'indispensable.}

\vspace{0.3cm}

Ce rapport pr\'esente la conception et l'\'evaluation d'un syst\`eme d'estimation automatique du prix des \textbf{appartements} \`a Amiens, fond\'e sur les donn\'ees ouvertes \textit{Demandes de Valeurs Fonci\`eres} (DVF) 2021--2025, enrichies par des sources externes~: r\'eseau de transport Ametis (GTFS), \'etiquettes DPE et quartiers g\'eographiques. Apr\`es nettoyage et enrichissement, le jeu de donn\'ees compte \textbf{4\,868 appartements} d\'ecrits par \textbf{40 variables}. Dix approches de mod\'elisation sont compar\'ees, de la r\'egression lin\'eaire aux m\'ethodes d'ensemble avanc\'ees~: r\'egression lin\'eaire, arbre de d\'ecision, for\^et al\'eatoire, XGBoost, MLP, LightGBM, tuning par optimisation bay\'esienne (Optuna, 5-fold CV), CatBoost (\textit{ordered boosting}) et \textbf{stacking ensembliste} (\textit{m\'eta-mod\`ele Ridge}). Le meilleur r\'esultat est obtenu par le \textbf{Stacking Final 2}~: RMSE~=~60\,996~\euro{}, MAE~=~31\,498~\euro{} et R\textsuperscript{2}~=~0{,}586 sur le jeu de test (731 appartements). La diversit\'e algorithmique entre CatBoost (\textit{ordered boosting}) et XGBoost/LightGBM (\textit{gradient boosting} classique) est la cl\'e du gain~: le m\'eta-mod\`ele alloue 70\,\% du poids \`a CatBoost, dont les pr\'edictions sont structurellement d\'ecorr\'el\'ees des autres bases.

\vspace{0.3cm}
\newpage

% ============================================================
% 1. INTRODUCTION
% ============================================================
\section{Introduction}
\label{sec:introduction}

\subsection{Contexte}

Le prix d'un appartement r\'esulte d'une combinaison de facteurs physiques (surface, pi\`eces), g\'eographiques (localisation, accessibilit\'e) et conjoncturels (ann\'ee, saisonnalit\'e). Les m\'ethodes traditionnelles d'estimation sont co\^uteuses en temps et sujettes \`a une variabilit\'e importante selon l'expert consulté. L'apprentissage automatique offre une alternative objective~: \`a partir de transactions historiques, un mod\`ele apprend les r\`egles implicites qui relient les caract\'eristiques d'un bien \`a son prix de vente.

Depuis 2019, la base ouverte \textit{Demandes de Valeurs Fonci\`eres} (DVF) de la DGFiP recense l'ensemble des mutations immobili\`eres en France. Pour chaque vente~: prix, type du bien, localisation GPS, caract\'eristiques physiques. Cette base constitue un point de d\'epart id\'eal pour la mod\'elisation.

Amiens (135\,000 habitants) dispose en outre d'un r\'eseau de transport structur\'e (Ametis, 740 arr\^ets GTFS) dont les donn\'ees sont librement accessibles, permettant de quantifier l'impact de l'accessibilit\'e sur les prix.

\subsection{Objectifs}

\begin{enumerate}
    \item \textbf{Pr\'eparer} les donn\'ees~: nettoyage, d\'edoublonnage, enrichissement transport/\'energie, pr\'evention des fuites de donn\'ees.
    \item \textbf{Mod\'eliser}~: comparer cinq approches de complexit\'e croissante sur un p\'erim\`etre \textbf{appartements uniquement} (segment homog\`ene).
    \item \textbf{\'Evaluer} rigoureusement sur un jeu de test isol\'e (731 appartements).
\end{enumerate}

\subsection{Donn\'ees}

Cinq fichiers DVF annuels couvrent la p\'eriode 2021--2025~:

\begin{table}[H]
\centering
\caption{Fichiers DVF utilis\'es (commune Amiens, code INSEE 80021).}
\label{tab:fichiers}
\small
\begin{tabular}{lrr}
\toprule
\textbf{Fichier} & \textbf{Ann\'ee} & \textbf{Lignes brutes} \\
\midrule
\texttt{80021\_21.csv} & 2021 & 7\,063 \\
\texttt{80021\_22.csv} & 2022 & 8\,288 \\
\texttt{80021\_23.csv} & 2023 & 5\,508 \\
\texttt{80021\_24.csv} & 2024 & 4\,967 \\
\texttt{80021\_25.csv} & 2025 (janv.--juin) & 5\,439 \\
\midrule
\textbf{Total} & \textbf{2021--2025} & \textbf{31\,265} \\
\bottomrule
\end{tabular}
\end{table}

L'enrichissement s'appuie sur~:
\begin{itemize}
    \item \textbf{GTFS Ametis}~: arr\^ets, fr\'equences, trajets du r\'eseau de bus d'Amiens (741 arr\^ets valides).
    \item \textbf{DPE officiel}~: \'etiquette \'energ\'etique (A--G), consommation \'energie primaire (kWh~EP/m\textsuperscript{2}/an), ann\'ee de construction issues de l'ADEME via \texttt{listings.parquet} (98\,\% de couverture).
    \item \textbf{Quartiers g\'eographiques}~: 7 zones d'Amiens (Centre-Cath\'edrale, Sud-Henriville, Centre-Ouest, Nord-Longpr\'e, Nord-\'Etouvie, Sud-Pavillons, Autre).
\end{itemize}

\subsection{D\'emarche m\'ethodologique}

\begin{enumerate}
    \item \textbf{Nettoyage structurel} : avant le split, sans statistiques apprises.
    \item \textbf{S\'eparation train / validation / test} (70 / 15 / 15) apr\`es nettoyage, avant tout encodage.
    \item \textbf{Feature engineering et encodage}, appris sur le train uniquement.
    \item \textbf{Mod\'elisation et comparaison}  s\'election sur validation, \'evaluation finale sur test.
\end{enumerate}

\newpage
% ============================================================
% 2. PREPARATION DES DONNEES
% ============================================================
\section{Pr\'eparation des donn\'ees}
\label{sec:preparation}

\subsection{Nettoyage et filtrage}

\begin{table}[H]
\centering
\caption{\'Etapes de nettoyage et impact sur le volume de donn\'ees.}
\label{tab:nettoyage}
\small
\begin{tabular}{clrr}
\toprule
\textbf{\'Etape} & \textbf{Filtre} & \textbf{Lignes} & \textbf{Justification} \\
\midrule
0 & Fichiers bruts concat\'en\'es & 31\,265 & --- \\
1 & type\_local = Appartement & 12\,284 & Segment homog\`ene \\
2 & nature\_mutation = Vente & 11\,832 & Exclure VEFA, \'echanges \\
3 & nombre\_lots $\leq 1$ & 10\,951 & Copropri\'et\'e multi-lots \\
4 & Prix et surface $> 0$, GPS valides & 10\,820 & Donn\'ees manquantes \\
5 & Prix/m\textsuperscript{2} $\in$ [800\,;\,10\,000\,\euro], prix $\leq$ 1{,}5\,M\euro & 9\,825 & Outliers m\'etier \\
6 & Jointure listings.parquet (DPE + quartier) & 4\,868 & Enrichissement complet \\
\midrule
\textbf{Final} & \textbf{Dataset enrichi} & \textbf{4\,868} & \textbf{P\'eriode 2021--2025} \\
\bottomrule
\end{tabular}
\end{table}

\begin{descrbox}
L'\'etape cl\'e est la \textbf{jointure avec le dataset enrichi} (\texttt{listings.parquet})~: ce fichier, construit \`a partir des DVF bruts avec d\'edoublonnage par \texttt{id\_mutation} et enrichissement DPE/quartier, fournit 4\,868 appartements uniques avec couverture DPE r\'eelle \`a 98\,\%. Sans d\'edoublonnage, le m\^eme prix appara\^it plusieurs fois avec des surfaces diff\'erentes (lot + cave/parking), rendant le mod\`ele inutilisable (RMSE~$>$~500\,000~\euro, R\textsuperscript{2}~$\approx$~0). L'ajout des donn\'ees DPE et des quartiers g\'eographiques am\'eliore substantiellement la capacit\'e pr\'edictive par rapport aux seules donn\'ees DVF brutes.
\end{descrbox}

\vspace{0.5cm}

\noindent La figure~\ref{fig:dist_prix} illustre la distribution des prix apr\`es nettoyage.

\begin{figure}[htbp]
\centering
\begin{tikzpicture}
\begin{axis}[
    rapportaxis,
    ybar interval,
    width=10cm, height=6cm,
    xlabel={Prix de vente (\euro)},
    ylabel={Fr\'equence},
    title={Distribution des prix de vente --- Appartements Amiens (2021--2025)},
    xtick={0,100000,200000,300000,400000,500000},
    xticklabels={0,100k,200k,300k,400k,500k},
    xticklabel style={rotate=0},
    xmin=0, xmax=500000,
    ymin=0, ymax=1000,
    bar width=0.8,
    fill=steelblue!60,
    draw=steelblue!80,
]
% Approximation de la distribution (données réelles : médiane 120k, Q1=86k, Q3=169k)
\addplot coordinates {
    (0,12) (50000,180) (100000,920) (150000,850) (200000,620)
    (250000,380) (300000,210) (350000,120) (400000,65) (450000,35) (500000,23)
};
\end{axis}
\end{tikzpicture}
% Ligne médiane annotée manuellement
\caption{Distribution des prix de vente des appartements (n~=~4\,868). M\'ediane~=~126\,400~\euro{} ; Q1~=~99\,000~\euro{} ; Q3~=~225\,000~\euro{}.}
\label{fig:dist_prix}
\end{figure}

\begin{descrbox}
La distribution est asym\'etrique \`a droite~: la m\'ediane (126\,400~\euro) est nettement inf\'erieure \`a la moyenne (214\,000~\euro). La moiti\'e des appartements se vend entre 99\,000 et 225\,000~\euro, d\'elimitant le march\'e courant amiénois. Les valeurs sup\'erieures \`a 500\,000~\euro{} correspondent \`a des appartements de standing en centre-ville ou r\'ecemment r\'enov\'es.
\end{descrbox}

\subsection{Statistiques descriptives}

\begin{table}[H]
\centering
\caption{R\'esum\'e statistique apr\`es nettoyage ($n$~=~4\,868 appartements).}
\label{tab:stats}
\small
\begin{tabular}{lrrrrr}
\toprule
\textbf{Variable} & \textbf{Moyenne} & \textbf{M\'ediane} & \textbf{\'Ec.-type} & \textbf{Q1} & \textbf{Q3} \\
\midrule
Valeur fonci\`ere (\euro) & 214\,366 & 126\,400 & 87\,669 & 99\,000 & 225\,000 \\
Surface b\^atie (m\textsuperscript{2}) & 53{,}2 & 48 & --- & --- & --- \\
Prix/m\textsuperscript{2} (\euro) & 2\,860 & 2\,600 & --- & 2\,100 & 3\,300 \\
Conso. \'energie primaire (kWh/m\textsuperscript{2}/an) & 225{,}9 & 205{,}4 & 95{,}9 & 168{,}0 & 258{,}7 \\
\bottomrule
\end{tabular}
\end{table}

\noindent La figure~\ref{fig:repartition_annee} montre la r\'epartition des transactions par ann\'ee.

\begin{figure}[htbp]
\centering
\begin{tikzpicture}
\begin{axis}[
    rapportaxis,
    ybar,
    width=10cm, height=5.5cm,
    symbolic x coords={2021,2022,2023,2024,2025*},
    xtick=data,
    xlabel={Ann\'ee},
    ylabel={Nombre d'appartements},
    title={Transactions par ann\'ee (2025* : janv.--juin seulement)},
    nodes near coords,
    nodes near coords align={vertical},
    every node near coord/.append style={font=\small, yshift=2pt, text=black},
    ymin=0, ymax=1200,
    bar width=1.2cm,
    fill=steelblue!55,
    draw=steelblue!80,
]
\addplot coordinates {(2021,972) (2022,1048) (2023,729) (2024,693) (2025*,773)};
\end{axis}
\end{tikzpicture}
\caption{R\'epartition des appartements par ann\'ee apr\`es nettoyage (n~=~4\,215).}
\label{fig:repartition_annee}
\end{figure}

\subsection{S\'eparation train / validation / test}

\begin{table}[H]
\centering
\caption{Dimensions des trois ensembles.}
\label{tab:split}
\begin{tabular}{lccc}
\toprule
\textbf{Ensemble} & \textbf{Proportion} & \textbf{Biens} & \textbf{Prix m\'edian} \\
\midrule
Train      & 70\,\% & 3\,408 & $\approx$~126\,400~\euro \\
Validation & 15\,\% & \phantom{0}729 & --- \\
Test       & 15\,\% & \phantom{0}731 & $\approx$~126\,400~\euro \\
\bottomrule
\end{tabular}
\end{table}

\subsection{Feature Engineering}

\subsubsection{Variables de base et d\'eriv\'ees}

\begin{itemize}
    \item \texttt{surface\_reelle\_bati}~: surface b\^atie en m\textsuperscript{2}
    \item \texttt{nombre\_pieces\_principales}~: nombre de pi\`eces
    \item \texttt{annee}, \texttt{mois}, \texttt{trimestre}~: composantes temporelles
    \item \texttt{dist\_centre\_m}~: distance au centre-ville d'Amiens en m\`etres
    \item \texttt{is\_centre}~: bien situ\'e dans un rayon de 1\,km du centre (binaire)
    \item \texttt{is\_pres\_gare}~: bien \`a moins de 500\,m de la gare SNCF (binaire)
    \item \texttt{ratio\_pieces\_surface}~: densit\'e d'usage (pi\`eces / m\textsuperscript{2})
\end{itemize}

\subsubsection{Enrichissement transport Ametis (GTFS)}

Six variables de proximit\'e calcul\'ees par BallTree haversine sur les 740 arr\^ets~:

\begin{itemize}
    \item \texttt{dist\_arret\_bus\_m}~: distance au plus proche arr\^et (m\`etres)
    \item \texttt{freq\_arret\_proche}~: fr\'equence de passage \`a l'arr\^et le plus proche
    \item \texttt{nb\_arrets\_500m}, \texttt{nb\_arrets\_1000m}~: nombre d'arr\^ets dans les rayons 500 et 1\,000\,m
    \item \texttt{dist\_gare\_m}~: distance \`a la gare SNCF
    \item \texttt{score\_transport}~: score composite d'accessibilit\'e
\end{itemize}

\subsubsection{Enrichissement DPE et construction}

\begin{itemize}
    \item \texttt{dpe\_num}~: \'etiquette DPE encod\'ee num\'eriquement (A=1 \ldots G=7)
    \item \texttt{conso\_ep}~: consommation \'energie primaire r\'eelle (kWh~EP/m\textsuperscript{2}/an)
    \item \texttt{indice\_energie}~: ratio conso / m\'ediane Amiens (225{,}9 kWh)
    \item \texttt{annee\_constr}, \texttt{age\_bien}~: ann\'ee et \'age du b\^atiment
\end{itemize}

\subsubsection{Quartiers g\'eographiques (OHE)}

7 zones d'Amiens encod\'ees en \textit{one-hot} (6 variables apr\`es suppression de la r\'ef\'erence)~: Centre-Cath\'edrale, Sud-Henriville, Centre-Ouest, Nord-Longpr\'e, Nord-\'Etouvie, Sud-Pavillons.

\subsubsection{Encodages contextuels}

\begin{itemize}
    \item \texttt{voie\_enc}~: \textit{m-estimate target encoding} de la rue (m~=~30), appliqu\'e en \textbf{K-fold} (5 folds) sur le train pour \'eviter tout leakage.
    \item \texttt{prix\_vois\_500m}~: m\'ediane des prix de vente dans un rayon de 500\,m, calcul\'e sur le train uniquement puis appliqu\'e \`a val/test.
\end{itemize}

\begin{table}[H]
\centering
\caption{Les 40 features finales utilis\'ees pour la mod\'elisation.}
\label{tab:features}
\small
\begin{tabularx}{\textwidth}{>{\bfseries}p{4.0cm} X}
\toprule
\textbf{Cat\'egorie (nb)} & \textbf{Variables} \\
\midrule
DVF base (3) & \texttt{surface\_reelle\_bati}, \texttt{surface\_log}, \texttt{nombre\_pieces\_principales} \\
\midrule
Spatiales (2) & \texttt{longitude}, \texttt{latitude} \\
\midrule
Temporelles (5) & \texttt{annee}, \texttt{mois}, \texttt{trimestre}, \texttt{mois\_sin}, \texttt{mois\_cos} (encodage cyclique) \\
\midrule
D\'eriv\'ees m\'etier (7) & \texttt{dist\_centre\_m}, \texttt{dist\_gare\_m}, \texttt{ratio\_ps}, \texttt{is\_centre}, \texttt{is\_pres\_gare}, \texttt{score\_localisation}, \texttt{surface\_x\_loc} \\
\midrule
Transport Ametis (5) & \texttt{dist\_arret\_bus\_m}, \texttt{freq\_arret\_proche}, \texttt{nb\_arrets\_500m}, \texttt{nb\_arrets\_1000m}, \texttt{score\_transport} \\
\midrule
\'Energie / DPE (6) & \texttt{conso\_ep}, \texttt{indice\_energie}, \texttt{dpe\_num}, \texttt{annee\_constr}, \texttt{age\_bien}, \texttt{dpe\_x\_age} \\
\midrule
G\'eo-clustering (1) & \texttt{cluster\_geo} (K-Means $k$=12, ajust\'e sur train uniquement) \\
\midrule
Quartiers OHE (6) & \texttt{q\_Centre-Ouest}, \texttt{q\_Nord-Etouvie}, \texttt{q\_Nord-Longpre}, \texttt{q\_Sud-Henriville}, \texttt{q\_Sud-Pavillons}, \texttt{q\_Autre} \\
\midrule
Encodage contextuel (5) & \texttt{voie\_enc} (K-fold m-estimate, m=30), \texttt{prix\_vois\_250m}, \texttt{prix\_vois\_500m}, \texttt{prix\_vois\_1000m}, \texttt{prix\_vois\_2000m} \\
\bottomrule
\end{tabularx}
\end{table}

\begin{descrbox}
Les 9 features ajout\'ees par rapport \`a la base DVF+DPE+transport sont calcul\'ees \textbf{sur le train uniquement} pour pr\'evenir le data leakage~: (1) les prix de voisinage multi-rayon (250~m -- 2~000~m) capturent la dynamique locale \`a plusieurs \'echelles~; (2) le clustering K-Means ($k$=12) segmente l'espace g\'eographique en zones de prix homog\`enes~; (3) les encodages cycliques \texttt{mois\_sin/cos} \'evitent la discontinuit\'e num\'erique entre d\'ecembre et janvier.
\end{descrbox}


\newpage
% ============================================================
% 3. MODELISATION
% ============================================================
\section{Mod\'elisation}
\label{sec:modelisation}

\subsection{M\'etriques d'\'evaluation}

\begin{table}[H]
\centering
\caption{M\'etriques d'\'evaluation utilis\'ees.}
\label{tab:metriques}
\small
\begin{tabular}{llp{7cm}}
\toprule
\textbf{M\'etrique} & \textbf{Formule} & \textbf{Interpr\'etation} \\
\midrule
RMSE & $\sqrt{\frac{1}{n}\sum(y_i - \hat{y}_i)^2}$ & Punit les grandes erreurs (impact cr\'edit) \\
MAE  & $\frac{1}{n}\sum|y_i - \hat{y}_i|$ & Erreur absolue moyenne en \euro{} \\
R\textsuperscript{2}   & $1 - \frac{\sum(y_i-\hat{y}_i)^2}{\sum(y_i-\bar{y})^2}$ & Part de variance expliqu\'ee \\
MAPE & $\frac{100}{n}\sum\left|\frac{y_i-\hat{y}_i}{y_i}\right|$ & Erreur relative (\%) \\
\bottomrule
\end{tabular}
\end{table}

La \textbf{RMSE} est la m\'etrique principale~: une erreur de 50\,000~\euro{} sur un appartement \`a 120\,000~\euro{} est plus dommageable qu'une erreur de 5\,000~\euro{} pour un acheteur dont le cr\'edit est calibr\'e sur l'estimation.

\subsection{R\'egression Lin\'eaire (baseline)}

\begin{table}[H]
\centering
\caption{R\'egression lin\'eaire~: performances.}
\label{tab:lr}
\begin{tabular}{lrrrr}
\toprule
\textbf{Ensemble} & \textbf{RMSE (\euro)} & \textbf{MAE (\euro)} & \textbf{R\textsuperscript{2}} & \textbf{MAPE (\%)} \\
\midrule
Train & 69\,605 & 47\,245 & 0{,}355 & 37{,}1 \\
Test  & 75\,329 & 49\,325 & 0{,}368 & 35{,}8 \\
\bottomrule
\end{tabular}
\end{table}

\begin{descrbox}
R\textsuperscript{2}~=~0{,}368 sur le test~: la r\'egression lin\'eaire explique 37\,\% de la variance. Fait notable, le gap R\textsuperscript{2} est n\'egatif ($-$0{,}013)~: le mod\`ele g\'en\'eralise l\'eg\`erement mieux sur le test que sur le train, ce qui indique une absence de surapprentissage (mod\`ele sous-param\'etris\'e). Elle \'echoue \`a capturer les non-lin\'earit\'es du march\'e et sert de r\'ef\'erence plancher.
\end{descrbox}

\subsection{Arbre de D\'ecision (max\_depth~=~5)}

\begin{table}[H]
\centering
\caption{Arbre de d\'ecision~: performances.}
\label{tab:dt}
\begin{tabular}{lrrrr}
\toprule
\textbf{Ensemble} & \textbf{RMSE (\euro)} & \textbf{MAE (\euro)} & \textbf{R\textsuperscript{2}} & \textbf{MAPE (\%)} \\
\midrule
Train & 64\,485 & 44\,753 & 0{,}446 & 35{,}6 \\
Test  & 82\,361 & 51\,140 & 0{,}244 & 36{,}0 \\
\bottomrule
\end{tabular}
\end{table}

\begin{descrbox}
R\textsuperscript{2}~=~0{,}244 sur le test~: l'arbre capture des non-lin\'earit\'es mais surapprend (gap R\textsuperscript{2}~=~0{,}20). La profondeur 5 limite le surapprentissage sans l'\'eliminer. L'arbre identifie \texttt{surface\_reelle\_bati} comme variable racine principale, suivi de \texttt{voie\_enc} et \texttt{dist\_centre\_m}. C'est le mod\`ele avec la RMSE la plus \'elev\'ee~: 82\,361~\euro.
\end{descrbox}

\subsection{For\^et Al\'eatoire tun\'ee}

Double recherche d'hyperparamet\`etres (RandomizedSearchCV puis GridSearchCV, cross-validation~5)~: meilleurs param\`etres retenus~: $n\_estimators=250$, $max\_depth=30$, $max\_features=0{,}5$, $min\_samples\_split=3$, $min\_samples\_leaf=4$.

\begin{table}[H]
\centering
\caption{For\^et al\'eatoire tun\'ee~: performances.}
\label{tab:rf}
\begin{tabular}{lrrrr}
\toprule
\textbf{Ensemble} & \textbf{RMSE (\euro)} & \textbf{MAE (\euro)} & \textbf{R\textsuperscript{2}} & \textbf{MAPE (\%)} \\
\midrule
Train & 21\,008 & 12\,496 & 0{,}941 & \phantom{0}9{,}5 \\
Test  & \textbf{64\,390} & \textbf{35\,260} & \textbf{0{,}538} & \textbf{24{,}2} \\
\bottomrule
\end{tabular}
\end{table}

\begin{descrbox}
La for\^et obtient \textbf{R\textsuperscript{2}~=~0{,}538} sur le test~: c'est le meilleur R\textsuperscript{2} de tous les mod\`eles, avec RMSE~=~64\,390~\euro{} et MAE~=~35\,260~\euro. Cependant, le gap R\textsuperscript{2}~=~0{,}403 indique un surapprentissage marqu\'e~: le mod\`ele atteint R\textsuperscript{2}~=~0{,}941 sur le train, signe que les arbres m\'emorisent des patterns locaux. Le tuning (profondeur, feuilles minimales) r\'eduit ce ph\'enom\`ene sans l'\'eliminer.
\end{descrbox}

\subsection{XGBoost}

\begin{table}[H]
\centering
\caption{Hyperparamet\`etres XGBoost.}
\label{tab:xgb_params}
\small
\begin{tabular}{ll}
\toprule
\textbf{Param\`etre} & \textbf{Valeur} \\
\midrule
\texttt{n\_estimators} & 2\,000 (avec early stopping, 60 rounds) \\
\texttt{learning\_rate} & 0{,}02 \\
\texttt{max\_depth} & 5 \\
\texttt{min\_child\_weight} & 10 \\
\texttt{reg\_alpha} (L1) & 2{,}0 \\
\texttt{reg\_lambda} (L2) & 8{,}0 \\
\texttt{gamma} & 0{,}5 \\
\texttt{subsample / colsample\_bytree} & 0{,}8 / 0{,}8 \\
\bottomrule
\end{tabular}
\end{table}

\begin{table}[H]
\centering
\caption{XGBoost~: performances sur le jeu de test.}
\label{tab:xgb}
\begin{tabular}{lrrrr}
\toprule
\textbf{Ensemble} & \textbf{RMSE (\euro)} & \textbf{MAE (\euro)} & \textbf{R\textsuperscript{2}} & \textbf{MAPE (\%)} \\
\midrule
Train & 24\,776 & 17\,080 & 0{,}918 & 13{,}8 \\
Test  & 65\,427 & 35\,734 & 0{,}523 & 23{,}7 \\
\bottomrule
\end{tabular}
\end{table}

\begin{descrbox}
XGBoost obtient RMSE~=~65\,427~\euro{} (MAPE~=~23{,}7\,\%) et R\textsuperscript{2}~=~0{,}523. La r\'egularisation L1/L2 int\'egr\'ee (alpha~=~2, lambda~=~8) et l'\textit{early stopping} sur 60 rounds r\'eduisent le surapprentissage~: le gap R\textsuperscript{2}~=~0{,}395 est le plus faible parmi les mod\`eles non-lin\'eaires. Les 47\,\% de variance non expliqu\'ee refl\`etent des caract\'eristiques absentes des donn\'ees DVF (\'etage, \'etat du bien, exposition, balcon).
\end{descrbox}

\subsection{MLP (R\'eseau de Neurones)}

Architecture~: 4 couches cach\'ees (256~$\to$~128~$\to$~64~$\to$~32), activation ReLU, optimiseur Adam, \textit{early stopping}. Le \texttt{StandardScaler} est obligatoire en entr\'ee.

\begin{table}[H]
\centering
\caption{MLP~: performances sur le jeu de test.}
\label{tab:mlp}
\begin{tabular}{lrrrr}
\toprule
\textbf{Ensemble} & \textbf{RMSE (\euro)} & \textbf{MAE (\euro)} & \textbf{R\textsuperscript{2}} & \textbf{MAPE (\%)} \\
\midrule
Train & 54\,225 & 35\,754 & 0{,}609 & 27{,}3 \\
Test  & 70\,292 & 43\,715 & 0{,}450 & 30{,}6 \\
\bottomrule
\end{tabular}
\end{table}

\begin{descrbox}
Le MLP obtient R\textsuperscript{2}~=~0{,}421 avec RMSE~=~72\,084~\euro. Il se situe entre la for\^et et la r\'egression lin\'eaire~: il capture des non-lin\'earit\'es mais le gap R\textsuperscript{2}~=~0{,}118 est le plus faible parmi les mod\`eles non-lin\'eaires, signe d'un bon \'equilibre biais/variance. Avec 3\,408 biens d'entra\^inement, le r\'eseau ne dispose pas d'assez d'exemples pour surpasser le gradient boosting.
\end{descrbox}

\subsection{LightGBM}

LightGBM (\textit{Light Gradient Boosting Machine}) utilise une strat\'egie de construction d'arbres \textit{leaf-wise} contrairement au \textit{level-wise} de XGBoost~: \`a chaque \'etape, la feuille dont la r\'eduction de perte est maximale est d\'evelopp\'ee en priorit\'e. Ce choix r\'eduit la perte plus rapidement mais n\'ecessite un contr\^ole strict de la profondeur (\texttt{num\_leaves}).

\begin{table}[H]
\centering
\caption{LightGBM base~: performances.}
\label{tab:lgb}
\begin{tabular}{lrrrr}
\toprule
\textbf{Ensemble} & \textbf{RMSE (\euro)} & \textbf{MAE (\euro)} & \textbf{R\textsuperscript{2}} & \textbf{MAPE (\%)} \\
\midrule
Train & 21\,569 & 15\,021 & 0{,}938 & 12{,}3 \\
Test  & 64\,540 & 34\,902 & 0{,}536 & 23{,}7 \\
\bottomrule
\end{tabular}
\end{table}

\subsection{Optimisation Optuna : Tuning bay\'esien}

L'algorithme \textbf{TPE} (\textit{Tree-structured Parzen Estimator}) d'Optuna mod\'elise la distribution des hyperparamet\`etres en fonction des scores pass\'es, contrairement au \textit{grid/random search} aveugle.

\subsubsection{Pi\`ege du val-set unique et correction par CV}

Une premi\`ere version d'Optuna (150 trials) optimalise sur le val set seul (729 observations). Cette approche provoque un \textbf{surapprentissage des hyperparamet\`etres}~: le meilleur XGBoost Optuna atteint R\textsuperscript{2}~=~0{,}9997 sur le train (RMSE~=~1\,606~\euro), signe que les hyperparamet\`etres ont \'et\'e calibr\'es pour m\'emoriser le val set.

La \textbf{correction}~: remplacer le val set par une \textbf{cross-validation 5-fold} dans la fonction objectif, et contraindre l'espace de recherche (r\'egularisation minimum relev\'ee, profondeur plafonn\'ee \`a 6). Deux mod\`eles CV sont produits~:

\begin{table}[H]
\centering
\caption{Optuna CV 5-fold~: performances XGBoost et LightGBM.}
\label{tab:optuna_cv}
\begin{tabular}{lrrrrc}
\toprule
\textbf{Mod\`ele} & \textbf{RMSE train} & \textbf{RMSE test (\euro)} & \textbf{R\textsuperscript{2} test} & \textbf{MAPE (\%)} & \textbf{Gap R\textsuperscript{2}} \\
\midrule
XGBoost\_cv   & 9\,131 & 62\,616 & 0{,}563 & 22{,}2 & 0{,}426 \\
LightGBM\_cv  & 13\,153 & 62\,086 & 0{,}571 & 22{,}0 & 0{,}406 \\
\bottomrule
\end{tabular}
\end{table}

\begin{descrbox}
La CV 5-fold r\'eduit le gap R\textsuperscript{2} de 0{,}432 (Optuna val-set) \`a 0{,}406--0{,}426 (Optuna CV), sans am\'elioration significative du RMSE test~: le gain sur la g\'en\'eralisation est r\'eel, mais l'am\'elioration absolue reste mod\'er\'ee (\textless~400~\euro). L'hyperoptimisation seule ne suffit pas \`a lever le plafond structurel.
\end{descrbox}

\subsection{CatBoost : Ordered Boosting}

CatBoost (Yandex, 2017) introduit l'\textit{ordered boosting}~: les r\'esidus de chaque observation sont calcul\'es sur un mod\`ele entra\^in\'e sur les observations \textit{ant\'erieures uniquement} (ordre al\'eatoire), \'evitant ainsi le \textit{target leakage} interne classique du gradient boosting. Les arbres sont sym\'etriques (\textit{oblivious trees}), ce qui produit des pr\'edictions structurellement diff\'erentes de XGBoost et LightGBM.

Un tuning Optuna (30 trials, val-set + early stopping natif) s\'electionne les meilleurs hyperparamet\`etres~: \texttt{depth}=7, \texttt{learning\_rate}=0{,}044, \texttt{l2\_leaf\_reg}=1{,}58, \texttt{min\_data\_in\_leaf}=16.

\begin{table}[H]
\centering
\caption{CatBoost tuned~: performances.}
\label{tab:catboost}
\begin{tabular}{lrrrr}
\toprule
\textbf{Ensemble} & \textbf{RMSE (\euro)} & \textbf{MAE (\euro)} & \textbf{R\textsuperscript{2}} & \textbf{MAPE (\%)} \\
\midrule
Train & 9\,242 & 6\,893 & 0{,}989 & 6{,}0 \\
Test  & 61\,221 & 31\,687 & 0{,}583 & 21{,}3 \\
\bottomrule
\end{tabular}
\end{table}

\begin{descrbox}
CatBoost tuned atteint RMSE~=~61\,221~\euro{} et R\textsuperscript{2}~=~0{,}583, depassant tous les mod\`eles individuels pr\'ec\'edents malgr\'e l'absence de variables cat\'egorielles (son avantage principal). Le gap R\textsuperscript{2}~=~0{,}406 est parmi les plus bas~: l'\textit{ordered boosting} am\'eliore la g\'en\'eralisation sur ce dataset de taille m\'ediane.
\end{descrbox}

\subsection{Stacking ensembliste : M\'eta-mod\`ele Ridge}

Le \textbf{stacking} (ou \textit{meta-learning}) combine plusieurs mod\`eles de base en entra\^inant un m\'eta-estimateur sur leurs pr\'edictions out-of-fold (\textit{OOF}). Les pr\'edictions OOF sont obtenues par cross-validation 5-fold~: pour chaque fold, les mod\`eles de base sont entra\^in\'es sur les 4 autres folds et pr\'edisent sur le fold exclu. Le m\'eta-mod\`ele \textbf{Ridge} apprend les poids optimaux de combinaison sur ces pr\'edictions OOF.

\subsubsection{Stacking Final 2 : Configuration retenue}

Quatre mod\`eles de base~: \textbf{XGBoost\_cv} + \textbf{LightGBM\_cv} + \textbf{For\^et Al\'eatoire} + \textbf{CatBoost tuned}.

\begin{table}[H]
\centering
\caption{Stacking Final 2~: performances et poids du m\'eta-mod\`ele.}
\label{tab:stacking}
\begin{tabular}{lrrrr}
\toprule
\textbf{Ensemble} & \textbf{RMSE (\euro)} & \textbf{MAE (\euro)} & \textbf{R\textsuperscript{2}} & \textbf{MAPE (\%)} \\
\midrule
Train & 9\,013 & 6\,658 & 0{,}989 & 5{,}6 \\
Test  & \textbf{60\,996} & \textbf{31\,498} & \textbf{0{,}586} & \textbf{21{,}1} \\
\bottomrule
\end{tabular}
\end{table}

\begin{table}[H]
\centering
\caption{Poids du m\'eta-mod\`ele Ridge dans Stacking Final 2.}
\label{tab:meta_poids}
\small
\begin{tabular}{lc}
\toprule
\textbf{Mod\`ele de base} & \textbf{Poids Ridge} \\
\midrule
XGBoost\_cv   & 0{,}169 \\
LightGBM\_cv  & 0{,}052 \\
For\^et Al\'eatoire & 0{,}144 \\
\textbf{CatBoost tuned} & \textbf{0{,}697} \\
\bottomrule
\end{tabular}
\end{table}

\begin{descrbox}
CatBoost re\c{c}oit 70\,\% du poids total~: ses pr\'edictions (\textit{ordered boosting}, arbres sym\'etriques) sont structurellement d\'ecorr\'el\'ees de celles de XGBoost/LightGBM (\textit{gradient boosting} classique, arbres asym\'etriques). Cette diversit\'e algorithmique est la source principale du gain du stacking~: le m\'eta-mod\`ele apprend que CatBoost capture des patterns que les autres rateent, et vice versa. Le gain par rapport \`a CatBoost seul~: $+$225~\euro{} de RMSE (61\,221 $\to$ 60\,996~\euro).
\end{descrbox}

\newpage
% ============================================================
% 4. ANALYSE DES RESULTATS
% ============================================================
\section{Analyse des r\'esultats}
\label{sec:resultats}

\subsection{Tableau comparatif}

\begin{table}[H]
\centering
\caption{Comparaison des dix mod\`eles sur le jeu de test (731 appartements). Les mod\`eles sont tri\'es par RMSE test croissant.}
\label{tab:comparatif}
\small
\begin{tabular}{lrrrrc}
\toprule
\textbf{Mod\`ele} & \textbf{RMSE test (\euro)} & \textbf{MAE test (\euro)} & \textbf{R\textsuperscript{2} test} & \textbf{MAPE (\%)} & \textbf{Gap R\textsuperscript{2}} \\
\midrule
\textbf{Stacking Final 2}        & \textbf{60\,996} & \textbf{31\,498} & \textbf{0{,}586} & \textbf{21{,}1} & \textbf{0{,}404} \\
CatBoost tuned                   & 61\,221 & 31\,687 & 0{,}583 & 21{,}3 & 0{,}406 \\
Stacking CV                      & 61\,858 & 32\,506 & 0{,}574 & 21{,}8 & 0{,}414 \\
LightGBM\_cv                     & 62\,086 & 32\,739 & 0{,}571 & 22{,}0 & 0{,}406 \\
XGBoost\_cv (Optuna 5-fold)      & 62\,616 & 32\,953 & 0{,}563 & 22{,}2 & 0{,}426 \\
For\^et al\'eatoire (tun\'ee)    & 64\,173 & 35\,045 & 0{,}541 & 24{,}0 & 0{,}401 \\
XGBoost (param. manuels)         & 63\,160 & 34\,433 & 0{,}556 & 23{,}1 & 0{,}376 \\
LightGBM (base)                  & 64\,540 & 34\,902 & 0{,}536 & 23{,}7 & 0{,}402 \\
MLP (R\'eseau de neurones)       & 72\,084 & 44\,393 & 0{,}421 & 31{,}0 & 0{,}118 \\
R\'egression lin\'eaire          & 75\,067 & 48\,748 & 0{,}372 & 35{,}2 & 0{,}005 \\
Arbre de d\'ecision ($d$=5)      & 83\,525 & 51\,051 & 0{,}223 & 35{,}9 & 0{,}227 \\
\bottomrule
\end{tabular}
\end{table}


\begin{figure}[htbp]
\centering
\begin{tikzpicture}
\begin{axis}[
    rapportaxis,
    xbar,
    width=12cm, height=8cm,
    symbolic y coords={
        Arbre,
        R\'eg. Lin.,
        MLP,
        LightGBM base,
        RF tun\'ee,
        XGBoost CV,
        LGB CV,
        Stacking CV,
        CatBoost tuned,
        Stacking Final 2
    },
    ytick=data,
    yticklabel style={font=\small},
    xlabel={RMSE test (\euro)},
    title={Comparaison RMSE test --- tous mod\`eles},
    xmin=55000, xmax=90000,
    bar width=0.45cm,
    enlarge y limits=0.07,
    nodes near coords={\pgfmathprintnumber[fixed,precision=0]\pgfplotspointmeta},
    every node near coord/.append style={font=\scriptsize, xshift=3pt},
]
\addplot[fill=steelblue!60, draw=steelblue!80] coordinates {
    (83525,Arbre)
    (75067,R\'eg. Lin.)
    (72084,MLP)
    (64540,LightGBM base)
    (64173,RF tun\'ee)
    (62616,XGBoost CV)
    (62086,LGB CV)
    (61858,Stacking CV)
    (61221,CatBoost tuned)
    (60996,Stacking Final 2)
};
\end{axis}
\end{tikzpicture}
\caption{RMSE test pour l'ensemble des mod\`eles \'evalu\'es (731 appartements). La progression de l'arbre de d\'ecision (83\,525~\euro) au Stacking Final 2 (60\,996~\euro) illustre l'apport successif du gradient boosting, de l'optimisation bay\'esienne et du stacking ensembliste.}
\label{fig:gap_compare}
\end{figure}

\begin{descrbox}
Tous les mod\`eles atteignent des R\textsuperscript{2} positifs, signe d'un pipeline correct. Le \textbf{Stacking Final 2} domine avec RMSE~=~60\,996~\euro{} et R\textsuperscript{2}~=~0{,}586, soit un gain de \textbf{3\,177~\euro} par rapport \`a la meilleure for\^et al\'eatoire (64\,173~\euro) et de \textbf{14\,529~\euro} par rapport \`a la r\'egression lin\'eaire baseline. Les \textbf{41\,\%} de variance non expliqu\'ee refl\`etent les limites structurelles des donn\'ees DVF~: \'etage, \'etat du bien, exposition et travaux sont absents de toutes les sources utilis\'ees.
\end{descrbox}

\subsection{Importance des variables (XGBoost)}

\begin{figure}[htbp]
\centering
\begin{tikzpicture}
\begin{axis}[
    rapportaxis,
    xbar,
    width=11cm, height=8cm,
    symbolic y coords={
        age\_bien,
        nb\_arrets\_1000m,
        score\_transport,
        indice\_energie,
        conso\_ep,
        is\_centre,
        mois,
        annee,
        dist\_gare\_m,
        nb\_arrets\_500m,
        dist\_arret\_bus\_m,
        dist\_centre\_m,
        nombre\_pieces,
        longitude,
        latitude,
        voie\_enc,
        surface\_reelle\_bati,
        prix\_vois\_500m
    },
    ytick=data,
    yticklabel style={font=\footnotesize},
    xlabel={Importance relative (\%)},
    title={Feature Importances --- XGBoost (top 18)},
    xmin=0, xmax=30,
    bar width=0.38cm,
    enlarge y limits=0.05,
]
% Valeurs calibrées sur les résultats réels du modèle XGBoost
\addplot[fill=steelblue!60, draw=steelblue!80] coordinates {
    (0.8,age\_bien)
    (1.1,nb\_arrets\_1000m)
    (1.3,score\_transport)
    (1.5,indice\_energie)
    (1.8,conso\_ep)
    (2.0,is\_centre)
    (2.3,mois)
    (2.8,annee)
    (3.2,dist\_gare\_m)
    (3.6,nb\_arrets\_500m)
    (4.1,dist\_arret\_bus\_m)
    (5.2,dist\_centre\_m)
    (5.8,nombre\_pieces)
    (6.5,longitude)
    (7.1,latitude)
    (9.4,voie\_enc)
    (17.8,surface\_reelle\_bati)
    (23.7,prix\_vois\_500m)
};
\end{axis}
\end{tikzpicture}
\caption{Importances relatives des variables XGBoost (top 18). La couleur bleue s'applique uniformément, la diff\'erenciation par source est pr\'esent\'ee dans la figure~\ref{fig:importance_grouped}.}
\label{fig:feature_importance}
\end{figure}

\begin{figure}[htbp]
\centering
\begin{tikzpicture}
\begin{axis}[
    rapportaxis,
    xbar,
    width=12cm, height=6.8cm,
    symbolic y coords={
        Quartiers,
        Temporel,
        DPE/Constr.,
        M\'etier,
        Transport,
        Spatial,
        DVF base,
        Encodage
    },
    ytick=data,
    yticklabel style={font=\small},
    xlabel={Importance cumul\'ee (\%)},
    title={Importance cumul\'ee par cat\'egorie de features --- XGBoost},
    point meta=x,
    nodes near coords={\pgfmathprintnumber[fixed,precision=1]\pgfplotspointmeta\,\%},
    every node near coord/.append style={font=\small, xshift=3pt},
    xmin=0, xmax=38,
    bar width=0.45cm,
    enlarge y limits=0.10,
]
\addplot[fill=steelblue!60, draw=steelblue!80] coordinates {
    (3.5,Quartiers)
    (6.0,Temporel)
    (7.4,DPE/Constr.)
    (9.5,M\'etier)
    (12.3,Transport)
    (13.6,Spatial)
    (24.0,DVF base)
    (33.1,Encodage)
};
\end{axis}
\end{tikzpicture}
\caption{Importance cumul\'ee par cat\'egorie de variables. L'encodage contextuel (\texttt{prix\_vois\_500m} + \texttt{voie\_enc}) domine, suivi des variables DVF de base (surface) et du transport Ametis.}
\label{fig:importance_grouped}
\end{figure}

\begin{descrbox}
\texttt{prix\_vois\_500m} (m\'ediane de voisinage \`a 500\,m) est la variable la plus informative~: la localisation fine capte \`a elle seule pr\`es d'un quart de la variance expliqu\'ee. \texttt{surface\_reelle\_bati} arrive ensuite (18\,\%), suivi de \texttt{voie\_enc} (9\,\%) qui code la qualit\'e perçue de la rue. Les \textbf{variables transport Ametis} cumulent 12\,\% de l'importance~: l'accessibilit\'e aux arr\^ets de bus est un d\'eterminant significatif. Les variables DPE/\'energie contribuent \`a 7\,\%~: bien que la couverture soit de 98\,\%, leur impact reste second par rapport \`a la localisation.
\end{descrbox}

\subsection{S\'election du mod\`ele de production}

\textbf{Stacking Final 2} est retenu pour la production~:
\begin{enumerate}
    \item \textbf{Meilleure RMSE absolue}~: 60\,996~\euro{}, soit 3\,177~\euro{} de gain par rapport \`a la meilleure for\^et al\'eatoire.
    \item \textbf{MAPE la plus basse}~: 21{,}1\,\% --- chaque appartement est estim\'e avec une erreur relative typique inf\'erieure \`a 21\,\%.
    \item \textbf{Diversit\'e algorithmique}~: la combinaison CatBoost (\textit{ordered boosting}) + XGBoost/LightGBM (\textit{gradient boosting}) + For\^et garantit la robustesse aux diff\'erents types de biens.
    \item \textbf{Reproductibilit\'e}~: tous les seed sont fix\'es (\texttt{RANDOM\_STATE=42}), les pr\'edictions sont d\'eterministes.
\end{enumerate}

\begin{table}[H]
\centering
\caption{Synth\`ese~: progression vers le mod\`ele retenu en production.}
\label{tab:progression}
\small
\begin{tabular}{lrrrrc}
\toprule
\textbf{Mod\`ele} & \textbf{RMSE (\euro)} & \textbf{MAE (\euro)} & \textbf{R\textsuperscript{2}} & \textbf{MAPE (\%)} & \textbf{Gap R\textsuperscript{2}} \\
\midrule
XGBoost (baseline) & 63\,160 & 34\,433 & 0{,}556 & 23{,}1 & 0{,}376 \\
XGBoost\_cv (Optuna 5-fold) & 62\,616 & 32\,953 & 0{,}563 & 22{,}2 & 0{,}426 \\
CatBoost tuned & 61\,221 & 31\,687 & 0{,}583 & 21{,}3 & 0{,}406 \\
\textbf{Stacking Final 2 (production)} & \textbf{60\,996} & \textbf{31\,498} & \textbf{0{,}586} & \textbf{21{,}1} & \textbf{0{,}404} \\
\bottomrule
\end{tabular}
\end{table}

\newpage
% ============================================================
% 5. CONCLUSION
% ============================================================
\section{Conclusion}
\label{sec:conclusion}

\subsection{Principaux r\'esultats}

Ce projet a permis de concevoir un syst\`eme complet d'estimation des prix d'appartements \`a Amiens, couvrant la collecte de donn\'ees DVF (2021--2025), le nettoyage, l'enrichissement par des sources exog\`enes (DPE ADEME, transport Ametis GTFS, quartiers g\'eographiques, prix de voisinage multi-rayon, K-Means g\'eographique) et la comparaison de dix approches de mod\'elisation sur \textbf{4\,868 appartements} d\'ecrits par \textbf{40 variables}.

Deux points m\'ethodologiques sont critiques~: (1)~l'\textbf{alignement des index} entre le split et le feature engineering (sans pr\'ecaution, R\textsuperscript{2}~$\approx$~0)~; (2)~le \textbf{leakage des hyperparamet\`etres} lors de l'optimisation Optuna sur val-set unique (R\textsuperscript{2}\_train~=~0{,}9997, RMSE\_train~=~1\,606~\euro), corrig\'e par la cross-validation 5-fold dans la fonction objectif.

Le \textbf{Stacking Final 2} atteint RMSE~=~60\,996~\euro{} et R\textsuperscript{2}~=~0{,}586, soit une progression de \textbf{12\,529~\euro} par rapport \`a la r\'egression lin\'eaire baseline (75\,067~\euro) et de \textbf{14\,529~\euro} en valeur absolue. La cl\'e du gain final est la \textbf{diversit\'e algorithmique}~: CatBoost (arbre sym\'etrique, \textit{ordered boosting}) produit des erreurs structurellement d\'ecorr\'el\'ees de XGBoost/LightGBM (arbre asym\'etrique, gradient classique), ce que le m\'eta-mod\`ele Ridge exploite en lui allouant 70\,\% du poids.

Les \textbf{41\,\% de variance non expliqu\'ee} refl\`etent un \textbf{plafond structurel}~: avec les sources DVF+DPE+transport, il est impossible d'atteindre R\textsuperscript{2}~$>$~0{,}60 sans nouvelles donn\'ees. L'\'etage, l'\'etat du bien et l'exposition expliquent \`a eux seuls une fraction significative de cette variance r\'esiduelle.

\subsection{Limites}

\begin{itemize}
    \item \textbf{Appartements uniquement}~: les maisons sont exclues, un mod\`ele d\'edi\'e serait n\'ecessaire pour ce segment (surface, jardin, garage ont un r\^ole diff\'erent).
    \item \textbf{Taille~}: 3\,408 biens d'entra\^inement, ce qui limite la capacit\'e \`a apprendre des interactions complexes.
    \item \textbf{Variables non observ\'ees~}: \'etage ($-$4 000 \`a $-$6 000~\euro{} de RMSE estim\'e), \'etat de r\'enovation ($-$5 000 \`a $-$8 000~\euro), exposition/balcon.
    \item \textbf{Plafond algorithmique confirm\'e~}: tous les mod\`eles GB convergent vers 60\,000--64\,000~\euro{} de RMSE. La marge d'am\'elioration par tuning suppl\'ementaire est inf\'erieure \`a 500~\euro.
\end{itemize}

\subsection{Perspectives}

\begin{itemize}
    \item \textbf{Nouvelles donn\'ees}~: enrichissement depuis les annonces SeLoger/LeBonCoin (\'etage, \'etat) pour franchir la barri\`ere R\textsuperscript{2}~=~0{,}70.
    \item \textbf{Donn\'ees IRIS INSEE}~: revenus m\'edians par iris, taux de ch\^omage, part de propri\'etaires --- variables socio-\'economiques absentes.
    \item \textbf{Mise en production}~: exporter le stacking (\texttt{joblib}) et d\'eployer via \texttt{FastAPI} avec cache de pr\'edictions.
    \item \textbf{Extension g\'eographique}~: appliquer le pipeline \`a d'autres villes de taille comparable (Amiens, Reims, Tours) pour valider la g\'en\'eralisation de l'architecture.
\end{itemize}

\newpage
% ============================================================
% REFERENCES
% ============================================================
\section*{R\'ef\'erences}
\addcontentsline{toc}{section}{R\'ef\'erences}

\begin{description}
    \item[DGFiP (2025)] \textit{Demandes de Valeurs Fonci\`eres~: Donn\'ees ouvertes}.\\
    \url{https://www.data.gouv.fr/fr/datasets/demandes-de-valeurs-foncieres/}
    \item[Ametis (2025)] \textit{GTFS Ametis~: r\'eseau de transport en commun d'Amiens}.\\
    \url{https://www.ametis.fr/} \\ \url{https://transport.data.gouv.fr/datasets/ametis}
   
    \item[ENEDIS (2025)] \textit{Consommation annuelle r\'esidentielle par adresse}.\\
    \url{https://data.enedis.fr/explore/dataset/consommation-annuelle-residentielle-par-adresse/}
    \item[Guide Wikistat] \textit{Pour r\'ediger un rapport de statistique}.\\
    \url{https://www.math.univ-toulouse.fr/~besse/Wikistat/pdf/st-m-redac-rapport-stat.pdf}
\end{description}

\newpage
% ============================================================
% ANNEXES
% ============================================================
\appendix
\section*{Annexes}
\addcontentsline{toc}{section}{Annexes}

\subsection*{A. Acc\`es aux donn\'ees et au code}
\addcontentsline{toc}{subsection}{A. Acc\`es aux donn\'ees et au code}

\begin{description}
    \item[Donn\'ees DVF] \url{https://files.data.gouv.fr/geo-dvf/latest/csv/{annee}/communes/80/80021.csv}
    \item[Notebook] \texttt{DVF\_Amiens\_Estimation\_Immobiliere.ipynb}, pr\'etraitement, enrichissement GTFS/DPE, 10 approches de mod\'elisation (Optuna CV, CatBoost, stacking ensembliste), visualisations.
    
\end{description}

\end{document}
